\subsection{问题三源码：往返场及配对对照}
\label{sec:appendix-q3}

两束截断与完整往返的机制比较要固定数据、尺度、光学参数及响应，只切换返回阶数；若分别重估参数，所得差值还包含估参变化。两束截断保留首轮返回光，完整往返则把后续每轮带着衰减和相位的返回光依次相加；逐次往返厚度对照据此组织两种场，再分别核对硅的全量厚度和碳化硅的沿用结果。

图\ref{fig:roundtrip}说明固定界面关系时返回场如何叠加，图\ref{fig:16}则展示分别拟合后的预测与残差。前者对应场的构造，后者对应整套估计方案；共同测试区段只固定了比较对象，不能单独隔离返回阶数的作用。

清单\ref{lst:q3-complete}中的配对求解让两模型共享粗候选、起点和评估上限，均重估全部光学参数及各角响应。分别估参的厚度差应与固定同一光学参数的纯场差分开，不能由其中一项代替另一项。

点集核对改变了最初的全量计算方案：每角$480$点的计划折分曾同时承担全量估计，但逐行核对发现主窗口实际包含每角$5392$个原始点。因此将计划折分保留为区段稳定性检验，并另用窗口内全部原始坐标和值重新估计厚度，避免以抽取点集冒称全量。
% src:交接/实验记录.json:559.类别;559.尝试;559.现象;559.决定
% src:求解/问题3/结果/解读核验_回炉3_仲裁后.json:防泄漏.每角抽样点数;防泄漏.全量每角点数

原始全量拟合给出的硅厚度为$2.17\,\mu\mathrm{m}$，输入保留每角$5392$个观测点。图\ref{fig:20}将这一全量点估计与折间结果分开标识；硅的厚度计算独立于碳化硅是否需要修正，必要物理条件的判断不取消硅的数值回答。
% src:求解/问题3/结果/回炉轮3候选/仲裁闭环/20260910_131501_81623_建模公式/厚度结果.json:核心指标.硅全量完整往返厚度_um
% src:求解/问题3/结果/解读核验_回炉3_仲裁后.json:防泄漏.全量每角点数

折间厚度分别为$1.73$和$16.04\,\mu\mathrm{m}$，各自对应既定连续区段的估计对象。它们描述区段选择带来的分支移动，既不平均成全量答案，也不把两端点当作全量厚度的概率区间。全量与折间分别保留之后，材料处理才有明确的比较基准。
% src:求解/问题3/结果/回炉轮3候选/仲裁闭环/20260910_131501_81623_建模公式/厚度结果.json:核心指标.硅折1完整往返条件厚度_um;核心指标.硅折2完整往返条件厚度_um

材料处理还须区分反射率预测与厚度修正：完整往返反射率区间的测试覆盖率为$73.91\%$，低于名义$90\%$。图\ref{fig:17}保留各材料、角度及区段的改善方向，表\ref{tab:08}则并列谱形误差与厚度；覆盖率的计量对象始终是反射率观测点，而修正厚度还需要相干性、空间重叠与仪器分辨率的观测依据。
% src:求解/问题3/结果/回炉轮3候选/仲裁闭环/20260910_131501_81623_建模公式/厚度结果.json:预测区间.4.经验覆盖率;预测区间.4.平均区间宽度_比例;预测区间.5.经验覆盖率;预测区间.5.平均区间宽度_比例;预测区间.6.经验覆盖率;预测区间.6.平均区间宽度_比例;预测区间.7.经验覆盖率;预测区间.7.平均区间宽度_比例;预测区间.20.经验覆盖率;预测区间.20.平均区间宽度_比例;预测区间.21.经验覆盖率;预测区间.21.平均区间宽度_比例;预测区间.22.经验覆盖率;预测区间.22.平均区间宽度_比例;预测区间.23.经验覆盖率;预测区间.23.平均区间宽度_比例;预测区间.36.经验覆盖率;预测区间.36.平均区间宽度_比例;预测区间.37.经验覆盖率;预测区间.37.平均区间宽度_比例;预测区间.38.经验覆盖率;预测区间.38.平均区间宽度_比例;预测区间.39.经验覆盖率;预测区间.39.平均区间宽度_比例;预测区间.52.经验覆盖率;预测区间.52.平均区间宽度_比例;预测区间.53.经验覆盖率;预测区间.53.平均区间宽度_比例;预测区间.54.经验覆盖率;预测区间.54.平均区间宽度_比例;预测区间.55.经验覆盖率;预测区间.55.平均区间宽度_比例;预测区间.4.名义覆盖率
% src:求解/问题3/结果/回炉轮3候选/仲裁闭环/20260910_131501_81623_建模公式/厚度结果.json:逐块评分.4.测试点数;逐块评分.5.测试点数;逐块评分.6.测试点数;逐块评分.7.测试点数;逐块评分.20.测试点数;逐块评分.21.测试点数;逐块评分.22.测试点数;逐块评分.23.测试点数;逐块评分.36.测试点数;逐块评分.37.测试点数;逐块评分.38.测试点数;逐块评分.39.测试点数;逐块评分.52.测试点数;逐块评分.53.测试点数;逐块评分.54.测试点数;逐块评分.55.测试点数

碳化硅沿用第二问的$10.80\,\mu\mathrm{m}$及$0.64$--$18.11\,\mu\mathrm{m}$光学情景范围，因为附件缺少相干、空间重叠和仪器分辨率资料，谱形改善尚未构成更换厚度的物理依据。表\ref{tab:q3-source-objects}把全量、折间和沿用关系分列，完整往返候选与实际采用厚度各有对应位置。
% src:求解/问题3/结果/回炉轮3候选/仲裁闭环/20260910_131501_81623_建模公式/厚度结果.json:核心指标.碳化硅正式基准厚度_um;核心指标.碳化硅正式条件范围下限_um;核心指标.碳化硅正式条件范围上限_um
% src:求解/问题3/结果/条件判定.json:必要条件;碳化硅条件分支

\begingroup
\small
\renewcommand{\arraystretch}{1.06}
\setlength{\LTpre}{0pt}
\setlength{\LTpost}{0pt}
\begin{longtable}{>{\centering\arraybackslash}p{0.28\textwidth}>{\centering\arraybackslash}p{0.25\textwidth}>{\centering\arraybackslash}p{0.39\textwidth}}
\caption{两材料全量厚度与折间结果分列}\label{tab:q3-source-objects}\\
\toprule
计算对象 & 厚度$/\mu\mathrm{m}$ & 数据或沿用关系\\
\midrule
\endfirsthead
\toprule
计算对象 & 厚度$/\mu\mathrm{m}$ & 数据或沿用关系\\
\midrule
\endhead
\bottomrule
\endfoot
硅全量 & $2.17$ & 每角$5392$个原始点\\
硅连续区段 & $1.73$、$16.04$ & 每角$480$点的计划切分\\
碳化硅同窗对照 & $2.7870\to2.8301$ & 两束$\to$完整场，采用值不变\\
碳化硅采用值 & $10.80$ & 沿用第二问同片基准\\
碳化硅情景范围 & $0.64$--$18.11$ & 沿用第二问光学情景\\
共同波数窗口 & --- & $1200$--$3800\,\mathrm{cm}^{-1}$\\
\end{longtable}
\endgroup
% src:求解/问题3/结果/回炉轮3候选/仲裁闭环/20260910_131501_81623_建模公式/厚度结果.json:案例.6.最优.厚度_um;案例.7.最优.厚度_um
% src:求解/问题3/结果/解读核验_回炉3_仲裁后.json:防泄漏.全量每角点数;防泄漏.每角抽样点数
% src:求解/问题3/结果/回炉轮3候选/仲裁闭环/20260910_131501_81623_建模公式/厚度结果.json:核心指标.硅全量完整往返厚度_um;核心指标.硅折1完整往返条件厚度_um;核心指标.硅折2完整往返条件厚度_um;核心指标.碳化硅正式基准厚度_um;核心指标.碳化硅正式条件范围下限_um;核心指标.碳化硅正式条件范围上限_um
% src:交接/结果声明_问题3.json:口径说明.硅全量完整往返厚度_um

沿用关系使碳化硅的处理前后保持同一个对象，而不是把换模型后的候选直接冠以修正厚度。现有观测对应保留基准的分支，两束与完整往返的候选仍成对保存；补充仪器资料后，可以沿同一数据对象重新判断高阶返回是否足以改变厚度。

已知光学参数下，匹配情景的最大厚度相对误差绝对值为$0.0521\%$；表\ref{tab:q3-known-thickness}按原情景列出全部$12$例直接厚度回收。低损耗、强损耗与无高阶负对照使用匹配的光学关系，色散失配组在估计时将色散设为零。真厚度、观测和生成条件保持不变，表中厚度单位均为$\mu\mathrm{m}$。
% src:求解/问题3/结果/公平对照.json:已知厚度验证.案例.3.相对误差_%;已知厚度验证.计划情景数;已知厚度验证.案例.0.口径;已知厚度验证.案例.9.口径

\begingroup
\small
\setlength{\LTpre}{0pt}
\setlength{\LTpost}{0pt}
\begin{longtable}{>{\centering\arraybackslash}p{0.20\textwidth}>{\centering\arraybackslash}p{0.16\textwidth}>{\centering\arraybackslash}p{0.17\textwidth}>{\centering\arraybackslash}p{0.17\textwidth}>{\centering\arraybackslash}p{0.14\textwidth}}
\caption{已知厚度的回收误差与包络包含}\label{tab:q3-known-thickness}\\
\toprule
生成情景 & 真厚度 & 估计厚度 & 误差幅度（\%） & 包络含真值\\
\midrule
\endfirsthead
\toprule
生成情景 & 真厚度 & 估计厚度 & 误差幅度（\%） & 包络含真值\\
\midrule
\endhead
低损耗 & $2.3500$ & $2.3501$ & $0.0048$ & 否\\
% src:求解/问题3/结果/公平对照.json:已知厚度验证.案例.0.情景;已知厚度验证.案例.0.真厚度_um;已知厚度验证.案例.0.估计厚度_um;已知厚度验证.案例.0.相对误差_%;已知厚度验证.案例.0.真厚度落入条件包络
低损耗 & $7.4500$ & $7.4498$ & $0.0023$ & 否\\
% src:求解/问题3/结果/公平对照.json:已知厚度验证.案例.1.情景;已知厚度验证.案例.1.真厚度_um;已知厚度验证.案例.1.估计厚度_um;已知厚度验证.案例.1.相对误差_%;已知厚度验证.案例.1.真厚度落入条件包络
低损耗 & $12.6500$ & $12.6500$ & $3.55\times10^{-5}$ & 否\\
% src:求解/问题3/结果/公平对照.json:已知厚度验证.案例.2.情景;已知厚度验证.案例.2.真厚度_um;已知厚度验证.案例.2.估计厚度_um;已知厚度验证.案例.2.相对误差_%;已知厚度验证.案例.2.真厚度落入条件包络
强损耗 & $2.3500$ & $2.3488$ & $0.0521$ & 是\\
% src:求解/问题3/结果/公平对照.json:已知厚度验证.案例.3.情景;已知厚度验证.案例.3.真厚度_um;已知厚度验证.案例.3.估计厚度_um;已知厚度验证.案例.3.相对误差_%;已知厚度验证.案例.3.真厚度落入条件包络
强损耗 & $7.4500$ & $7.4496$ & $0.0056$ & 否\\
% src:求解/问题3/结果/公平对照.json:已知厚度验证.案例.4.情景;已知厚度验证.案例.4.真厚度_um;已知厚度验证.案例.4.估计厚度_um;已知厚度验证.案例.4.相对误差_%;已知厚度验证.案例.4.真厚度落入条件包络
强损耗 & $12.6500$ & $12.6496$ & $0.0031$ & 否\\
% src:求解/问题3/结果/公平对照.json:已知厚度验证.案例.5.情景;已知厚度验证.案例.5.真厚度_um;已知厚度验证.案例.5.估计厚度_um;已知厚度验证.案例.5.相对误差_%;已知厚度验证.案例.5.真厚度落入条件包络
无高阶负对照 & $2.3500$ & $2.3501$ & $0.0029$ & 否\\
% src:求解/问题3/结果/公平对照.json:已知厚度验证.案例.6.情景;已知厚度验证.案例.6.真厚度_um;已知厚度验证.案例.6.估计厚度_um;已知厚度验证.案例.6.相对误差_%;已知厚度验证.案例.6.真厚度落入条件包络
无高阶负对照 & $7.4500$ & $7.4503$ & $0.0036$ & 否\\
% src:求解/问题3/结果/公平对照.json:已知厚度验证.案例.7.情景;已知厚度验证.案例.7.真厚度_um;已知厚度验证.案例.7.估计厚度_um;已知厚度验证.案例.7.相对误差_%;已知厚度验证.案例.7.真厚度落入条件包络
无高阶负对照 & $12.6500$ & $12.6502$ & $0.0019$ & 否\\
% src:求解/问题3/结果/公平对照.json:已知厚度验证.案例.8.情景;已知厚度验证.案例.8.真厚度_um;已知厚度验证.案例.8.估计厚度_um;已知厚度验证.案例.8.相对误差_%;已知厚度验证.案例.8.真厚度落入条件包络
色散失配 & $2.3500$ & $2.2629$ & $3.7082$ & 否\\
% src:求解/问题3/结果/公平对照.json:已知厚度验证.案例.9.情景;已知厚度验证.案例.9.真厚度_um;已知厚度验证.案例.9.估计厚度_um;已知厚度验证.案例.9.相对误差_%;已知厚度验证.案例.9.真厚度落入条件包络
色散失配 & $7.4500$ & $6.7823$ & $8.9622$ & 否\\
% src:求解/问题3/结果/公平对照.json:已知厚度验证.案例.10.情景;已知厚度验证.案例.10.真厚度_um;已知厚度验证.案例.10.估计厚度_um;已知厚度验证.案例.10.相对误差_%;已知厚度验证.案例.10.真厚度落入条件包络
色散失配 & $12.6500$ & $11.7205$ & $7.3479$ & 否\\
% src:求解/问题3/结果/公平对照.json:已知厚度验证.案例.11.情景;已知厚度验证.案例.11.真厚度_um;已知厚度验证.案例.11.估计厚度_um;已知厚度验证.案例.11.相对误差_%;已知厚度验证.案例.11.真厚度落入条件包络
\bottomrule
\end{longtable}
\endgroup

点估计回收误差与包络包含是不同问题：即使最优厚度贴近真值，已搜索到的近优分支仍可能集中在真值同一侧。这里以同一训练目标最优值的$10\%$为容差构成合成近优包络，包含判定使用未舍入端点；匹配组的$1/9$是模拟计数，不是实测置信水平。实测多情景包络还联合切分与光学扰动，其宽度来源与这组固定光学参数检验不同。
% src:求解/问题3/结果/公平对照.json:已知厚度验证.包络构造;已知厚度验证.条件包络经验包含率;已知厚度验证.完成匹配情景数;已知厚度验证.厚度覆盖解释

往返场实现见清单\ref{lst:q3-complete}，其来源与结果记录中的重算入口一致。程序依赖Python、NumPy及SciPy；工作簿的解压、数值单元格读取及散列核对由标准库完成。附件、数据档案和第二问基准的相对位置固定，输入索引见清单\ref{lst:q3-products}；不读取旧第三问程序的估计值作为新答案。
% src:求解/问题3/结果/优胜来源.json:结果重算入口;估计器版本

各角响应先用训练数据确定中心与尺度，再从物理反射率列中扣除低阶背景投影，使背景变化与条纹响应分开拟合。物理响应使用非负增益，并对增益与背景系数施加收缩。背景阶数、增益阶数和收缩强度均由训练内连续留段选择，两种场共用候选、起点与评估上限。校准数据构造反射率区间，测试数据用于评分；内层比较未完成时沿用预先固定的响应设置，不从部分候选中择优。朴素模型胜出只改变预测模型标签，硅厚度仍由完整往返场估计。
% src:求解/问题3/结果/公平对照.json:案例.0.训练内选择;案例.0.搜索状态;指标含义

一次重算得到的是独立保存的候选产物，不自动替换本文采用的声明。清单\ref{lst:q3-products}把估计、数值核对与汇总的现有步骤分别对应到产物；汇总程序按案例拆出逐块评分和扰动记录，并保留碳化硅的前问基准。随后形成的解读核验更新声明措辞，厚度仍对应原有核心指标。新候选需重新完成这些核对，不能沿用旧运行的比较结论。
% src:求解/问题3/结果/优胜来源.json:复核入口;结果重算入口;禁止混用
% src:求解/问题3/结果/回炉轮3候选/仲裁闭环/20260910_131501_81623_建模公式/厚度结果.json:核心指标

\begingroup
\renewcommand{\lstlistingname}{清单}
\lstset{basicstyle=\ttfamily\footnotesize,columns=fullflexible,
  keepspaces=true,showstringspaces=false,breaklines=true,
  numbers=none,frame=single,captionpos=t}
\begin{lstlisting}[caption={第三问的输入与结果对应关系},label={lst:q3-products}]
输入（相对工作根）：
数据/附件1.xlsx 至 数据/附件4.xlsx
交接/数据档案.json
求解/问题2/结果/基准交接.json

唯一估计入口：求解/问题3/升格3/求解.py
源码SHA256：
653c3f11046b42d388403f1342de2bae5d0e47733a8d1986290cba237b7ae45f
对应：求解/问题3/结果/公平对照.json 的 实现SHA256

原估计产物：求解/问题3/升格3/结果/
隔离重算产物：本次运行目录/结果/
公平对照.json；厚度结果.json；区间覆盖.json
结果声明_问题3.json（候选声明）；执行状态.json
训练内选择_<材料>_折<折号>.json
基础估计_<材料>_折<折号>.json
预测_<材料>_折<折号>.csv
配对重采样_<材料>.json；已知厚度验证.json
扩展计算的完成范围以实际记录为准。

既有核对：求解/问题3/裁决_升格.py
调用：求解/问题3/核验_升格标准库.py
产物：日志/升格裁决_问3核验/同口径比较.json
      日志/升格裁决_问3核验/独立公式与评分核验.json
既有汇总：求解/问题3/整理_升格裁决.py
产物：求解/问题3/结果/候选同口径比较.json
      求解/问题3/结果/裁决核验.json
      求解/问题3/结果/同口径对照.json
      求解/问题3/结果/灵敏度.json
      求解/问题3/结果/声明自检指标.json
      求解/问题3/结果/条件判定.json
      求解/问题3/结果/优胜来源.json
      交接/结果声明_问题3.json
本轮解读核对：求解/问题3/核验解读_回炉轮2.py
核对对象：求解/问题3/结果/回炉轮2候选/升格3复算/
核对产物：求解/问题3/结果/解读核验_回炉轮2.json
当前声明：交接/结果声明_问题3.json
\end{lstlisting}
\endgroup

上述核对、汇总程序会写入各自的结果、交接和日志目录，因此不在隔离重算命令中执行。清单\ref{lst:q3-run}只调用当前估计器，将其所有数值输出重定向到新建目录；既有汇总产物作为只读对应索引，不与本次候选混用。

\begingroup
\renewcommand{\lstlistingname}{清单}
\lstset{basicstyle=\ttfamily\footnotesize,columns=fullflexible,
  keywordstyle={},commentstyle={},stringstyle={},
  keepspaces=true,showstringspaces=false,breaklines=true,
  breakatwhitespace=false,numbers=left,numberstyle=\tiny,
  breakautoindent=true,breakindent=0pt,
  postbreak={\makebox[1em][l]{\textcolor{gray}{\ensuremath{\hookrightarrow}}}},
  numbersep=5pt,xleftmargin=1.5em,frame=single,
  captionpos=t,tabsize=4,aboveskip=0.8em,belowskip=0.8em}
\lstinputlisting[language=Python,firstline=1,lastline=674,caption={往返场与同段对照的完整实现},label={lst:q3-complete}]{../求解/问题3/升格3/求解.py}
\lstinputlisting[language=Python,firstline=675,lastline=900,firstnumber=675,title={\small 清单\ref{lst:q3-complete}（续）：扰动与验证}]{../求解/问题3/升格3/求解.py}
\lstinputlisting[language=Python,firstline=901,lastline=979,firstnumber=901,title={\small 清单\ref{lst:q3-complete}（续）：结果发布与记录}]{../求解/问题3/升格3/求解.py}
\lstinputlisting[language=Python,firstline=980,lastline=1098,firstnumber=980,title={\small 清单\ref{lst:q3-complete}（续）：执行入口}]{../求解/问题3/升格3/求解.py}
\endgroup

隔离复算从工作根执行，命令见清单\ref{lst:q3-run}。输入散列由估计器核对，源码散列与清单\ref{lst:q3-products}及原结果保存值逐一比对。程序按时间预算保存各阶段结果；需结合执行状态和每例搜索记录判断计算完成范围，进程正常结束本身不保证候选数值逐位相同。

\begingroup
\renewcommand{\lstlistingname}{清单}
\lstset{basicstyle=\ttfamily\footnotesize,columns=fullflexible,
  keywordstyle={},commentstyle={},stringstyle={},
  keepspaces=true,showstringspaces=false,breaklines=true,
  breakatwhitespace=false,numbers=none,frame=single,
  breakautoindent=true,breakindent=0pt,
  postbreak={\makebox[1em][l]{\textcolor{gray}{\ensuremath{\hookrightarrow}}}},
  captionpos=t,tabsize=4}
\begin{lstlisting}[language=bash,caption={独立保存数值与失败日志的复算命令},label={lst:q3-run}]
run_dir="$(mktemp -d 日志/问题三复算.XXXXXX)" || exit 1
printf '本次复算目录：%s\n' "$run_dir"
python3 -B - "$run_dir" >"$run_dir/运行日志.txt" 2>&1 <<'PYCODE'
import hashlib
import importlib.util
import json
import signal
import sys
from pathlib import Path

sys.dont_write_bytecode = True
source_path = Path("求解/问题3/升格3/求解.py").resolve()
expected_digest = (
    "653c3f11046b42d388403f1342de2bae5d0"
    "e47733a8d1986290cba237b7ae45f"
)
actual_digest = hashlib.sha256(source_path.read_bytes()).hexdigest()
recorded_digest = json.loads(
    Path("求解/问题3/结果/公平对照.json").read_text(encoding="utf-8")
)["实现SHA256"]
if actual_digest != expected_digest or actual_digest != recorded_digest:
    raise RuntimeError("源码摘要与本文采用的结果不一致")


def stop_at_budget(signum, frame):
    raise TimeoutError("达到二十分钟上限，保留已写出结果")


signal.signal(signal.SIGALRM, stop_at_budget)
signal.alarm(1200)
module_spec = importlib.util.spec_from_file_location(
    "question_three", source_path
)
model = importlib.util.module_from_spec(module_spec)
module_spec.loader.exec_module(model)
model.RESULT = Path(sys.argv[1]).resolve() / "结果"
try:
    model.main()
finally:
    signal.alarm(0)
PYCODE
\end{lstlisting}
\endgroup

复算结果由点集、场模型和材料处理共同定位，既能核对硅的全量答案，也能追踪碳化硅沿用基准的依据。补充材料据此展开参数来源索引与各路线比较，将数值对应回同一套观测条件。
